\honest{Availability}{Eliezer Yudkowsky}{2007}

The availability heuristic is judging how frequent or likely an event is by how easily examples of it come to mind. \nb{The idea and the name come from Tversky and Kahneman (1973).}

My evidence is a 1978 study by Lichtenstein and colleagues. Its subjects thought accidents killed about as many people as disease, and that homicide was more common than suicide. In fact disease kills about sixteen times as many people as accidents, and suicide is ``twice as frequent as homicide.''

A 1979 follow-up found that the errors correlated strongly, 0.85 and 0.89, with how often two newspapers reported each cause of death. I admit this does not show whether reporting shapes memory or vividness shapes both. Either way, ``an availability bias is at work.''

Selective reporting is a major cause, and it is worse now than for our ancestors, who had ``usually at most one layer of selective reporting'' between them and an event. I give no source for that. Today a report may pass through six bloggers before it reaches you, and far more happens in a far larger world, so far less of it reaches you: a stronger filter and a larger bias. You hear about Bill Gates far more than his numbers justify, and you may compare your success to his and feel worse. The reverse also holds: 19\% of the world lives on less than \$1 a day, and they do not write a fifth of the blog posts you read. \nb{That is a point about comparison and mood, not about judging how often something happens.}

Availability also ``seems'' to make people treat what has never happened as impossible. When there has been no recent flood, people refuse flood insurance even when it is heavily subsidized and priced far below its fair value. I quote Kunreuther and colleagues: men on flood plains are ``prisoners of their experience,'' and recent floods seem to set ``an upward bound'' on the losses they think worth worrying about. \nb{Yudkowsky's paper on global risks, from which much of this post is taken, gives these words to Kunreuther and colleagues in its 2006 draft and to Robert Kates (1962) in the version published in 2008.} Dams prevent small floods, so people take fewer precautions; floods become rarer but each does so much more damage that the average yearly damage rises.

From floods I draw a rule. The wise would reason from small hazards to the possibility of large ones; instead, experience of small hazards sets a ceiling on perceived risk. A society protected against minor hazards ``takes no action against major risks'' and builds on flood plains; a society used to minor hazards guards against those and not against the rare large one. \nb{The paragraph before had said ``reduced precautions,'' which is less than no action.} ``Memory is not always a good guide to probabilities in the past, let alone in the future.''
